\section{Related Work}\label{sec:related}

\sys{} builds on three lines of research spanning direct cache access, receive-ring and congestion management, and modern NIC-CPU interfaces.

\paragraph{Direct cache access and DDIO optimization.}
Direct Cache Access~\cite{huggahalli2005dca} and Intel DDIO~\cite{intelddio} steer inbound DMA into host LLC, yet Farshin et al.~\cite{farshin2019llc,farshin2020ddio} and ResQ~\cite{tootoonchian2018resq} demonstrated that oversized receive rings cause Leaky DMA and severe tail-latency degradation for co-located tenants. Wang et al.~\cite{wang2022ddio,wang2026sumeru} model cache-optimal inbound I/O, Yuan et al.~\cite{yuan2021ioalloc} redesign hardware LLC way allocation for I/O, PacketMill~\cite{farshin2021packetmill} pairs cache reservation with software metadata optimization, Shenango~\cite{ousterhout2019shenango} prefetches packet lines, and Sepia~\cite{song2026sepia} uses page coloring to eliminate DDIO set-conflict misses. These host-side techniques are complementary to \sys{}, but cannot absorb 16--128~MiB bursts larger than the host LLC without spilling to host DRAM or dropping packets.

\paragraph{Receive-ring sizing, host congestion control, and on-NIC buffering.}
ShRing~\cite{pismenny2023shring} shares one LLC-sized receive ring across cores (at the cost of locking or skew fallbacks), and rxBisect~\cite{pismenny2025rxbisect} removes locking by splitting allocation from delivery rings; both remain constrained by host memory conservation under bursts. HostCC~\cite{agarwal2023hostcc}, HostNet~\cite{vuppalapati2024hostnet}, Hostping~\cite{liu2023hostping}, and RHCC~\cite{wan2025rhcc} diagnose and reactively throttle intra-host I/O congestion via end-to-end CCA signals after host queues build. \ceio{}~\cite{liu2025ceio} introduced proactive NIC-side credit control and on-NIC elastic buffering on BlueField-3, achieving strong gains over HostCC and ShRing. \sys{} shares \ceio{}'s goal of proactive NIC-side burst absorption, but eliminates \ceio{}'s dual-path fast/slow split, per-flow RMT rule rewriting, phase-exclusivity pauses, 2P1C software merge ring, host DMA pull API, and $>$1K-flow churn cliff by unifying all traffic onto one credit-gated push FIFO.

\paragraph{Coherent, MMIO, streaming, and SmartNIC datapaths.}
CC-NIC~\cite{schuh2024ccnic} and CXL-NIC~\cite{cxlnic2025} explore cache-coherent NICs over UPI and CXL, while nanoPU~\cite{ibanez2021nanopu} and Dagger~\cite{lazarev2021dagger} tightly integrate packet queues into CPU registers or coherent FPGAs. Ens\={o}~\cite{sadok2023enso} replaces descriptors with contiguous streaming pipes on a custom FPGA, and recent work~\cite{liew2025hotos,liew2026asplos} studies non-coherent MMIO push for transmit. Wei et al.~\cite{wei2023smartnic} and Liu et al.~\cite{liu2026smartnic} characterize internal contention on SoC SmartNICs, and classic SPSC queue studies~\cite{lamport1983specifying,maffione2019spsc} optimize cache-aware producer-consumer rings. \sys{} applies cache-pinned, credit-gated one-way push directly to the receive path on unmodified non-coherent PCIe Gen5 silicon with full hardware IOMMU protection.
